\section{Introduction}
\label{sec:introduction}

The diagnosis of parasitic infections, particularly those caused by soil-transmitted helminths (STH), traditionally relies on manual microscopic examination of stool smears, such as the Kato-Katz technique. This manual process is labor-intensive, time-consuming, and highly dependent on the skill of the technician, often leading to inconsistencies and low sensitivity, especially in cases of light infection. The automation of this diagnostic workflow using computational methods presents a significant opportunity to improve efficiency, accuracy, and accessibility in global health settings.

Deep learning, specifically convolutional neural networks (CNNs), has emerged as a powerful tool for automating medical image analysis. While previous studies have demonstrated the viability of using object detectors like YOLO or Faster R-CNN for locating parasites \cite{viet2019parasite, wang2024automated}, a systematic comparison of modern architectures on this specific problem domain is lacking. Furthermore, the choice of architecture involves trade-offs between speed, accuracy, and the ability to capture detailed morphological information, such as the exact shape of an egg, which can be crucial for differentiating between species or from background artifacts \cite{cringoli2024automating}.

This paper presents a comprehensive benchmark of multiple state-of-the-art object detection and instance segmentation models for the task of identifying parasite eggs in microscopic images. We evaluate a diverse set of architectures, including one-stage detectors (RetinaNet, SSD, YOLO) and two-stage detectors (Faster R-CNN, Mask R-CNN), on a publicly available dataset. The primary objective is to quantitatively assess and compare their performance using the Mean Average Precision (mAP) metric, thereby providing insights into the most effective and practical models for automating parasite egg detection.

The remainder of this paper is structured as follows: Section \ref{sec:methodology} details the dataset, the models evaluated, and our experimental setup. Section \ref{sec:results} presents and discusses the comparative results of the benchmark. Finally, Section \ref{sec:conclusion} concludes our findings and suggests directions for future work.
